\documentclass[10pt,a4paper]{amsart}
\usepackage[margin=19mm]{geometry}
\usepackage{amsmath,amssymb,mathtools}
\usepackage[scr=rsfs,cal=euler]{mathalfa}
\usepackage{xfrac}
\usepackage[unicode,hidelinks,bookmarks=false]{hyperref}
\newcommand{\ie}{i.e.\ }
\newcommand{\cf}{cf.\ }
\newcommand{\resp}{respectively\ }
\newcommand{\Subsection}{\subsection}
\providecommand{\nobreakdash}{\nobreak\hbox{-}}
\setlength{\emergencystretch}{2em}
\multlinegap0pt
\usepackage{amssymb}
\newtheorem{proposition}{Proposition}
\title{Algorithm for direct sampling from conditional distributions of toric models}
\author{Shuhei Mano, Nobuki Takayama}
\date{Source: arXiv:2110.14992v3 (CC BY 4.0)}
\begin{document}
\maketitle

\noindent\emph{Source and licence.} Algorithm for direct sampling from conditional distributions of toric models. Version arXiv:2110.14992v3 (CC BY 4.0), distributed by arXiv under the Creative Commons Attribution 4.0 International licence, http://creativecommons.org/licenses/by/4.0/, as recorded in the arXiv licence metadata for this exact version and shown on the abstract page https://arxiv.org/abs/2110.14992v3. Full licence notice: \texttt{licenses/arxiv-2110.14992-CC-BY-4.0.txt}.\par\smallskip

\noindent\textit{Editorial scope.} Counted unit: the article's proposition prop:no-l-way (source0.tex lines 1559--1564). The article's complete original proof of it is delivered with it (source0.tex lines 1566--1612, one \texttt{proof} environment of 47 lines). No mathematics is written, repaired or re-derived: the statement and the proof are the article's own lines, in the article's own notation, and no retained line is substituted or re-flowed.  Retained uncounted as the source-local context the proof needs: lines 876--882.  Nothing was omitted from any retained span, so the package carries no omission record.  No retained line contains a citation, so the wrapper carries no bibliography and no bibliography entry.  No retained line contains a figure environment, an image-inclusion command or drawing code, so no image is delivered and the package declares no asset dependency.  Editorial formatting only, in addition: an AMS \texttt{amsart} preamble that restates the article's own declarations and macros used by the retained lines, namely the environments \texttt{proposition} on separate counters, together with the standard packages those lines need.  The retained environments print in this document as Proposition 1 in order of appearance, whereas the article prints the same items with the numbering of its own full document.  The complete native source of the article as distributed in the arXiv e-print for this exact version is bundled unchanged as \texttt{sources/117792/source0.tex}.
\begin{equation}
  {}_kF_{k-1}(c_1,\ldots,c_k;d_2,\ldots,d_k;z)
  :=\sum_{u\ge 0}\frac{(c_1)_u\cdots(c_k)_u}
  {(d_2)_u\cdots(d_k)_u}\frac{z^u}{u!}, \quad
  k\in\mathbb{N}\setminus\{1\},
  \label{chyp}
\end{equation}

\begin{proposition}\label{prop:no-l-way}
  The normalization constant of the conditional distribution
  of the binary no-$l(\ge 2)$-way interaction model is
  proportional to the classical univariate hypergeometric
  polynomial \eqref{chyp} with $k=2^{l-1}$.
\end{proposition}  

\begin{proof}
  The joint state space $\mathcal{I}_V=[2]^l$ can be identified
  with the hypercube and the sufficient statistics are the facets.
  For each state $i_V\in[2]^l$, there exists a shortest path to
  the state $(1\cdots 1)$ whose length is the number of coordinates
  of $i_V$ occupied by 2. Then, $u(i_V)$ can be expressed by
  an alternating sum of sufficient statistics and $u_{1\cdots 1}$.
  For example, the state $(1212)$ has a path
  $(1212)\to(1112)\to(1111)$, which can expressed as
  \[
    u_{1212}=u_{1\cdot12}-u_{1112}=u_{1\cdot12}-u_{111\cdot}
    +u_{1111}.
  \]
  For such an expression, $u(i_V)!$ can be represented by
  the Pochhammer symbol with subscript $u_{1\cdots 1}$, and
  $y_{i}^{u}$ can be represented by power of $y_{1\cdots1}$.
  For $u_{1212}!$,
  \[
    u_{1212}!
    =(u_{1\cdot12}-u_{111\cdot})!
    (u_{1\cdot12}-u_{111\cdot}+1)_{u_{1111}}
  \]
  and $y_{1212}^{u_{1212}}\propto y_{1212}^{u_{1111}}$.
  For the state $(1112)$, the unique shortest path to $(1111)$ gives
  \[
    u_{1112}!=(u_{111\cdot}-u_{1111})!=\frac{u_{111\cdot}!}
    {(-u_{111\cdot})_{u_{1111}}}(-1)^{u_{1111}}
  \]
  and $y_{1222}^{u_{1222}}\propto y_{1222}^{-u_{1111}}$.
  These observations lead to 
  \begin{equation}
    \frac{y^u}{u!}\propto
    \frac{\prod_{\{i:i\in[2]^l,~\text{number~of}~1~\text{is~odd}\}}
    (c_i)_{u_{1\cdots1}}}
    {\prod_{\{i:i\in[2]^l,~\text{number~of}~1~\text{is~even}\}}
    (d_i)_{u_{1\cdots1}}}
    \left(
    \frac{
    \prod_{\{i:i\in[2]^l,~\text{number~of}~1~\text{is~even}\}}y_i}
    {\prod_{\{i:i\in[2]^l,~\text{number~of}~1~\text{is~odd}\}}y_i}
    \right)^{u_{1\cdots1}},
    \label{toric2n}     
  \end{equation}
  which is the form of the summand of the classical univariate
  hypergeometric polynomial \eqref{chyp} with $k=2^{l-1}$ and
  $u=u_{1\cdots1}$. 
\end{proof}

\end{document}
